\chapter{K-nearest neighbors (k-NN)}
\label{chap:ch3-knn}

\section{Introduction to k-NN}
\label{sec:ch3-intro}

The \textbf{k-nearest neighbors (k-NN)} algorithm is a simple, yet powerful, method
used for classification and regression tasks in machine learning. The main idea
behind k-NN is to predict the output for a given input by looking at the 'k'
closest training examples in the feature space.

The k-NN algorithm works by finding the 'k' data points that are closest
to the query point in the feature space. The proximity between data points
is typically measured using distance metrics, such as \textbf{Euclidean distance} or
\textbf{Manhattan distance}. Once the closest neighbors are identified, the algo-
rithm makes predictions based on the majority class (in classification) or
the average value (in regression) of these neighbors.

\subsection{K-NN procedure}
\label{sec:ch3-knn-procedure}

The k-NN algorithm consists of the following steps:
\begin{enumerate}
  \item \textbf{Select the number of neighbors (k)}: Choose a value for 'k', which
    represents the number of nearest neighbors to consider for prediction.
  \item \textbf{Calculate distances}: Compute the distance between the query point
    and all data points in the training set.
  \item \textbf{Sort and select the nearest neighbors}: Sort the training data points
    based on the calculated distances and select the 'k' nearest neighbors.
  \item \textbf{Predict the output}:
    \begin{itemize}
      \item For \textbf{classification}, assign the class label that is most frequent
        among the 'k' neighbors.
      \item For \textbf{regression}, calculate the average of the target values of the 'k'
        neighbors and use it as the prediction.
    \end{itemize}
\end{enumerate}

\subsection{Choosing the Value of k}
\label{sec:ch3-choosing-k}

The value of 'k' is crucial to the performance of the k-NN algorithm. A small
value of 'k' can make the model sensitive to noise in the data, leading to
overfitting. A larger value of 'k' makes the model more robust but may lead
to underfitting if it averages out too many diverse neighbors. Typically, the
optimal value of 'k' is chosen through cross-validation.

\subsection{Advantages and Limitations}
\label{sec:ch3-advantages-limitations}

The k-NN algorithm has several advantages:
\begin{itemize}
  \item \textbf{Simplicity}: It is easy to understand and implement.
  \item \textbf{Non-parametric}: No assumptions about the underlying data distribution
    are made, making it versatile for various data types.
  \item \textbf{Lazy learning}: There is no explicit training phase. The model learns at
    the time of prediction.
\end{itemize}
However, it also has some limitations:
\begin{itemize}
  \item \textbf{Computationally expensive}: The algorithm requires calculating the dis-
    tance between the query point and all training points, which can be slow for
    large datasets.
  \item \textbf{Sensitive to the choice of k and distance metric}: The performance
    of the algorithm is highly dependent on the selection of 'k' and the distance
    metric used.
  \item \textbf{Memory-intensive}: Since the algorithm stores all training data, it requires
    significant memory, especially for large datasets.
\end{itemize}
Despite these limitations, the k-NN algorithm remains a fundamental and
widely used technique in machine learning for both classification and regression
problems.

\section{Worked-out Example}
\label{sec:ch3-worked-out}

Let's illustrate the k-NN algorithm with a simple classification example using
a small dataset.

\begin{figure}[h]\centering
% FIGURE PLACEHOLDER (uncomment & replace with \includegraphics when image is available):
% \fbox{\parbox{0.7\linewidth}{\centering\textit{[Figure from PDF page 27: 2D scatter plot with red points (class 0) at (1,2) and (2,4), blue points (class 1) at (3,3), (5,4), (6,5). A black query point at (4,4) with dashed lines showing distances $d=2$ to (2,4), $d=1$ to (5,4), and $d=1.41$ to (3,3). A small data table shows the coordinates and classes: (1,2)$\to$0, (2,4)$\to$0, (3,3)$\to$1, (5,4)$\to$1, (6,5)$\to$1.]}}}
\caption{Worked-out example for k-NN.}\label{fig:ch3-worked-out}
\end{figure}

\textbf{Dataset:} Consider the dataset on Figure~\ref{fig:ch3-worked-out} of points in a 2D space, where
each point is labeled as either \textbf{Class 0} or \textbf{Class 1}. We want to classify a new
point $(4, 4)$ based on this dataset using the k-NN algorithm.

\textbf{Step 1: Choose the value of k} For simplicity, let's choose $k = 3$. This
means we will consider the 3 closest neighbors to make our prediction.

\textbf{Step 2: Calculate the distances} We will calculate the Euclidean distance
between the new point $(4, 4)$ and each of the points in the dataset. The Eu-
clidean distance between two points $(x_1, y_1)$ and $(x_2, y_2)$ is given by:
\begin{equation}\label{eq:ch3-euclidean}
d = \sqrt{(x_2 - x_1)^2 + (y_2 - y_1)^2}.
\end{equation}
The distances are calculated as follows:

\begin{center}
\begin{tabular}{c|c|c}
\toprule
\textit{Point} & Distance to $(4,4)$ & Class \\
\midrule
$(1,2)$ & $\sqrt{(4-1)^2 + (4-2)^2} = \sqrt{9+4} = \sqrt{13} \approx 3.61$ & 0 \\
$(2,4)$ & $\sqrt{(4-2)^2 + (4-4)^2} = \sqrt{4} = 2$ & 0 \\
$(3,3)$ & $\sqrt{(4-3)^2 + (4-3)^2} = \sqrt{1+1} = \sqrt{2} \approx 1.41$ & 1 \\
$(5,4)$ & $\sqrt{(4-5)^2 + (4-4)^2} = \sqrt{1+0} = \sqrt{1} = 1$ & 1 \\
$(6,5)$ & $\sqrt{(4-6)^2 + (4-5)^2} = \sqrt{4+1} = \sqrt{5} \approx 2.24$ & 1 \\
\bottomrule
\end{tabular}
\end{center}

\textbf{Step 3: Sort and select the nearest neighbors:} The three points closest
to $(4,4)$ are: $(5,4)$, distance $= 1$; $(3,3)$, distance $= 1.41$; $(2,4)$, distance $= 2$.
These three points belong to the classes $1$, $1$ and $0$ respectively.

\textbf{Step 4: Prediction} Since the majority of the nearest neighbors (2 out of 3)
belong to \textbf{Class 1}, we predict that the point $(4, 4)$ belongs to \textbf{Class 1}. Thus,
the k-NN algorithm classifies the point $(4, 4)$ as \textbf{Class 1} based on the majority
class of its 3 nearest neighbors.

\section{Formal description of the k-NN procedure}
\label{sec:ch3-formal}

We consider a dataset of pairs $\cD = \{(X_i, Y_i), i = 1, \ldots, N\}$ where $X_i \in \cX$ and
$Y_i \in \cY$. Usually $\cX$ is a subset of $\R^d$ while $\cY$ can either be a discrete set $\cC$
(when there are several categories) or $\{0,1\}$ for binary classification or $\cY = \R$
for regression.

\begin{tcolorbox}[advancedbox]
\textbf{Advanced topics 3.1.} \textit{The coordinates of a point of $x \in \cX$ are called
``features'' and are supposed to be informative (predictive) on its cor-
responding $y \in \cY$ (category or value). As we will see latter, when $d$
is large the performance of the procedure deteriorates. For this reason,
when the objects live in a high dimensional space a first step is done
to extract the relevant ``features'' from the data. This will lower the
dimension and make the input more useful for the goal (regression or
categorization). PCA or similar techniques can also be invoked.}
\end{tcolorbox}

Take now $x \in \cX$ a new point to be predicted. We compute all distances
$\norm{X_i - x}$ for $i \leq N$ and obtain a permutation $\pi(x)$ of $\{1, \ldots, N\}$ that:
\begin{equation}\label{eq:ch3-permutation}
\underbrace{\norm{X_{\pi_1(x)} - x}}_{\text{closest to } x}
\leq \underbrace{\norm{X_{\pi_2(x)} - x}}_{\text{second closest to } x}
\leq \ldots \leq \underbrace{\norm{X_{\pi_N(x)} - x}}_{\text{furthest from } x}.
\end{equation}

Suppose that $k$ is a given integer, $k \geq 1$; this is the parameter of the
algorithm. For the binary regression the prediction is the following:
\begin{equation}\label{eq:ch3-regression}
\hat{f}_{\cD}(x) := \frac{1}{k} \sum_{i=1}^{k} Y_{\pi_i(x)}
\end{equation}

For the binary classification:
\begin{equation}\label{eq:ch3-classification}
\hat{f}_{\cD}(x) :=
\begin{cases}
1 & \text{if } \mathrm{card}\{i \leq k : Y_{\pi_i(x)} = 1\} \geq \mathrm{card}\{i \leq k : Y_{\pi_i(x)} = 0\} \\
0 & \text{otherwise}
\end{cases}
\end{equation}

where for a set $S$ we denote $\mathrm{card}(S)$ its cardinality. Note that when the number
of $Y_i$ (among the $k$ nearest neighbors) that belong to class $0$ and $1$ are equal
this algorithm biases favorably the class $1$. To avoid this bias either use an
odd value of $k$ or choose at random the ties.

For multi-class settings the function $\hat{f}_{\cD}(x)$ is analogously defined as re-
turning the class with cardinality larger than any other among the $k$-nearest
neighbors of $x$. Same remark as before applies about the possibility of ties.

\section{Error estimation with respect to Bayes predictor}
\label{sec:ch3-error-bayes}

In this section, we present Cover's proof \cite{turinici5} that when $k = 1$ the k-Nearest
Neighbors (k-NN) algorithm achieves an error rate that is less than twice the
Bayes error rate in the limit of large datasets.

Recall the notation~(2.1) of the law of $(X, Y)$ and introduce the probability
densities $f_1, \ldots, f_k$ such that an object in category $i$ gives rise to an observation
$x$ according to density $f_i$. Formally
\begin{equation}\label{eq:ch3-fi}
f_i(x) = \PP_{(X,Y) \sim \nu}[X = x \mid Y = i].
\end{equation}

Also introduce the (posterior) distributions:
\begin{equation}\label{eq:ch3-etai}
\eta_i(x) = \PP_{(X,Y) \sim \nu}[Y = i \mid X = x].
\end{equation}

To ease notation we will write $\E[\cdot]$ instead of $\E_{(X,Y) \sim \nu}[\cdot]$.

\begin{tcolorbox}[advancedbox]
\textbf{Advanced topics 3.2.} \textit{Note that if all functions $f_i$ have disjoint sup-
port the Bayes error will be zero and the following result is trivial and
not informative. This is the case when a given $x \in \cX$ corresponds to a
unique category $y \in \cY$.}
\end{tcolorbox}

We first need a preliminary result:

\begin{lemma}\label{lem:ch3-nearest}
Let $x$ be a fixed point and assume that with probability one a
$x \in \cX$ is either a continuity point of all $f_i$, $i \leq M$ or a point of non-zero
probability measure. Let $X_1, \ldots, X_n$ be i.i.d variables. Denote $Z_n$ the nearest
neighbor of $x$ from the collection $X_1, \ldots, X_n$. Then $Z_n \xrightarrow{n \to \infty} x$ with probability
one.
\end{lemma}

\begin{proof}
We accept the result which is somehow technical, see \cite[page 23]{turinici5} for
details.
\end{proof}

We will denote by $M$ the number of categories, $R^*$ the risk of the Bayes
estimator, $R_n$ the risk of the k-NN estimator for the i.i.d. variables $X_1, \ldots,$
$X_n$ and $R = \lim_{n \to \infty} R_n$.

\begin{theorem}\label{thm:ch3-cover}
We take $k = 1$ (one nearest neighbor). Assume that with
probability one a $x \in \cX$ is either a continuity point of all $f_i$, $i \leq M$ or a
point of non-zero probability measure. Then when the loss function is $l_{binary}$
in~(2.6):
\begin{equation}\label{eq:ch3-cover-bound}
R^* \leq R \leq R^* \left( 2 - \frac{M}{M-1} R^* \right).
\end{equation}
\end{theorem}

\begin{tcolorbox}[advancedbox]
\textbf{Advanced topics 3.5.} \textit{The theorem works with same proof for sepa-
rable metric spaces.}
\end{tcolorbox}

\begin{proof}
We will prove the assertion for $M = 2$ categories and refer to \cite[pages 25-26]{turinici5} for the general case. Denote $\eta(x) = \eta_1(x)$, cf. equation~\eqref{eq:ch3-etai}.
From~(2.12) we know that the Bayes risk is: $R^* = \E[\ind_{Y \neq f_{B_0}(X)}] =
\E[\min\{\eta(X), 1 - \eta(X)\}]$. On the other hand invoking Lemma~\ref{lem:ch3-nearest} the near-
est neighbor of $X$ tends to $X$ so by continuity hypothesis (the precise use
of the continuity requires some details that can be found in \cite[page 24]{turinici5})
the variable $Y_{\pi_1(X)}$ follows a Bernoulli distribution with parameter $\eta(X)$.
We obtain the risk
\begin{align}
R &= \E[\ind_{Y_{\pi_1(X)} \neq Y}] = \E\left[\E[\ind_{Y_{\pi_1(X)} \neq Y} \mid X]\right] = \E\left[\PP[Y_{\pi_1(X)} \neq Y \mid X]\right] \notag \\
  &= \E\left[\sum_{y \in \{0,1\}} \PP[Y = y, Y_{\pi_1(X)} = 1 - y \mid X]\right] \label{eq:ch3-risk-sum}
\end{align}
Note that conditional to $X$, variables $Y$ and $Y_{\pi_1(X)}$ are independent as
$\pi_1(X)$ depend on the given sample that does not involve $X$ and $Y$. We can
thus continue
\begin{align}
R &= \E\left[\sum_{y \in \{0,1\}} \PP[Y = y \mid X] \cdot \PP[Y_{\pi_1(X)} = 1 - y \mid X]\right] \notag \\
  &= \E\left[2 \eta(X)(1 - \eta(X))\right]. \label{eq:ch3-risk-final}
\end{align}
Denote $R(X) = \min\{\eta(X), 1 - \eta(X)\}$. Since $\E[2\eta(X)(1 - \eta(X))] =
\E[2R(X)(1 - R(X))]$ by concavity we obtain $R \leq 2\E[R(X)] \cdot (1 -
\E[R(X)]) = 2R^*(1 - R^*)$ hence the conclusion.
\end{proof}

\begin{remark}\label{rem:ch3-error-rate}
This proof shows that the k-NN algorithm's error rate is at
most twice the Bayes error, making it a robust non-parametric classifier.
As $k$ increases, the error rate will decrease.
\end{remark}

\section{Higher dimensions}
\label{sec:ch3-higher-dim}

\subsection{Datasets: dimension versus number of objects}
\label{sec:ch3-datasets}

The k-NN works in a special regime where the notion of `closeness' plays
a role. To understand it better we will need first some information on the
modern datasets. Each dataset $\cD$ is a collection of $n_{\cD}$ objects and each object
lives in a space with $d_{\cD}$ dimensions. We plot in figure~\ref{fig:ch3-datasets} this information
for some modern datasets. The most striking fact is that for recent datasets
such as those in 'NLP' and 'video' categories the number of objects is \textbf{smaller}
than the dimension of a single object while for historical datasets (the 'stat'
category) the number of objects is roughly two orders of magnitude higher
than the dimension of an object. This fact has far reaching implication that
will be described in the next sections.

\begin{figure}[h]\centering
% FIGURE PLACEHOLDER (uncomment & replace with \includegraphics when image is available):
% \fbox{\parbox{0.7\linewidth}{\centering\textit{[Figure from PDF page 32: Log-log scatter plot of Number of Objects in Dataset (y-axis, $10^0$ to $10^8$) versus Dimension of Object (x-axis, $10^0$ to $10^{12}$). Image datasets plotted: MNIST, CIFAR-10, SVHN, ImageNet, COCO, PASCAL VOC. Stat datasets: Iris, Wine, Boston Housing, Breast Cancer, Digits. NLP datasets: Wikipedia, IMDB Reviews, Enwik8, Common Crawl, OpenWebText. Video datasets: Kinetics-700, UCF101.]}}}
\caption{The number of objects versus the dimension of a single object
for historical datasets (label 'stat') and recent ones (labels 'image', 'video',
'NLP').}\label{fig:ch3-datasets}
\end{figure}

\subsection{Information contained in the distance}
\label{sec:ch3-info-distance}

Consider now the following question: suppose we take at random points in the
$d$ dimensional unit cube. As argued in the previous section taking $n$ points with
$n \sim d$ is maybe a good way to go. We compute the following metric; for each
point we take the distance to its nearest neighbor and rescale it by the median
of all distances between the all points in the set. We plot in figure~\ref{fig:ch3-min-distance} the
evolution of this normalized minimum distance with respect to the dimension.
It is seen that the minimum distance is informative for low dimensions i.e., is
different from the median distance, but as dimension increases the minimum
distance approaches the median distance so the notion of nearest neighbor
looses any meaning, all points are almost equally spaced. This means that
methods based on distances may work poorly in practice.

\begin{figure}[h]\centering
% FIGURE PLACEHOLDER (uncomment & replace with \includegraphics when image is available):
% \fbox{\parbox{0.7\linewidth}{\centering\textit{[Figure from PDF page 33: Line plot showing the (rescaled) minimum distance as a function of dimension (x-axis from $10^0$ to $10^3$, log scale; y-axis from 0.0 to 1.0). The min rescaled distance curve rises from about 0.2 at dimension 1 toward 1.0 at dimension $10^3$, with a shaded 99\% confidence interval band.]}}}
\caption{The distance to nearest neighbor (renormalized with respect to the
median distance of all points) as function of dimension. We take 1000 different
simulations, $N = 100$ points in each simulation and plot the empirical 99\%
confidence interval too.}\label{fig:ch3-min-distance}
\end{figure}

\begin{tcolorbox}[advancedbox]
\textbf{Advanced topics 3.7.} \textit{This can be formalized by computing, for a
$d$-dimensional normal sampling, the variance of the distance (squared)
between $d$ points and compare with the mean value of the distance.}
\end{tcolorbox}

\begin{tcolorbox}[advancedbox]
\textbf{Advanced topics 3.8.} \textit{On the contrary, distance to hyper-planes be-
haves well with increasing dimension; this remark motivates other suc-
cessful statistical approaches.}
\end{tcolorbox}

Not all is lost. It may still be that some dimensions are more important
than the other and a classification may still be possible, as illustrated in the
next section. For instance, when data is on a curve in a very high dimensional
space this the distance between the points is relevant again.

\subsection{Curse of dimensionality}
\label{sec:ch3-curse-dim}

We investigate here how the dimension of the data influences the quality of the
k-NN procedure. Consider a binary classification with data $x = (x_1, \ldots, x_d) \in
[-1/2, 1/2]^d$ following a uniform distribution on the cube and label $y$ equal to
$1$ if $x_1 \geq 0$ and $0$ otherwise. In this case there exists indeed a low dimensional
decision criterion but also dimensions that do not carry any information where
data is random.

The Bayes estimator is simply $\ind_{x_1 \geq 0}$ and has null error. Let us see how
k-NN performs when the dimension $d$ increases. Note that even when $d = 1$
the error is not always null but vanishes when the number $n$ of samples goes to
infinity. When $d$ is large $n$ will not be enough to sample well the whole cube.
So in this case the classification is possible but using distances is not a good
way to reach it. This will only work when there is low noise and data lies in a
low dimensional manifold.

\begin{figure}[h]\centering
% FIGURE PLACEHOLDER (uncomment & replace with \includegraphics when image is available):
% \fbox{\parbox{0.7\linewidth}{\centering\textit{[Figure from PDF page 34: Line plot of \% correct labels (y-axis from 50 to 100) versus dimension (x-axis, log scale from 1 to 200). Blue curve: k-NN (median) dropping from ~100\% at $d=1$ toward ~60\% at $d=200$, with shaded $Q_1$-$Q_3$ range. Red horizontal line: random choice at 50\%.]}}}
\caption{k-NN performance with respect to the dimension. We take in each
case $N = 100$ points and $300$ independent datasets.}\label{fig:ch3-curse-dim}
\end{figure}

\section{Computer lab : k-NN}
\label{sec:ch3-computer-lab}

\begin{tcolorbox}[remarkbox]
\textbf{Exercise 3.1} (accuracy vs dimension)\textbf{.} \textit{Test kNN for the Iris, MNIST
and CIFAR10 datasets (attention the last one is resource intensive).
You should obtain more than 90\% accuracy on first two but only a third
for the last one. Can you explain? Obtain some statistics involving the
classes of the nearest $k$ neighbors for objects in each dataset.}
\end{tcolorbox}
